% ============================================================================
\chapter{Memory Management}
% ============================================================================

\section{The Allocator Interface (\texttt{IAllocator})}

All allocators implement the pure-virtual interface:

\begin{lstlisting}[caption={IAllocator interface}]
class IAllocator {
public:
    virtual void*  alloc(usize size, usize alignment = 8) = 0;
    virtual void   free(void* ptr) = 0;
    virtual void   reset() = 0;
    virtual usize  usedMemory()      const = 0;
    virtual usize  totalSize()       const = 0;
    virtual usize  peakMemory()      const = 0;
    virtual usize  allocationCount() const = 0;
    virtual const char* name()       const = 0;
};
\end{lstlisting}

The alignment helper computes the number of padding bytes required to
bring a pointer to the next alignment boundary:

\begin{equation}
\text{padding}(p, a) =
\begin{cases}
  0 & \text{if } p \bmod a = 0 \\
  a - (p \bmod a) & \text{otherwise}
\end{cases}
\end{equation}

where $a$ must be a power of two. The fast bitwise form is
$(a - (p \mathbin{\&} (a-1))) \mathbin{\&} (a-1)$.


\paragraph{Invariant 4.1 --- Alignment}

  Every allocation returned by any \texttt{IAllocator} must satisfy:
  $(\text{address} \bmod \text{alignment}) = 0$.
  The engine assumes 8-byte minimum alignment by default.



\section{Linear Allocator}

\begin{definition}[Linear Allocator]
A linear (or ``bump-pointer'') allocator maintains a single cursor
$c$ into a contiguous buffer $[s, s+N)$. Allocation advances the cursor:
$c' = c + \text{padding}(c, a) + \text{size}$.
The only valid ``free'' operation is a full reset: $c \gets s$.
\end{definition}

\textbf{Complexity}: $O(1)$ allocation, $O(1)$ reset. No per-allocation
overhead beyond alignment padding.

\textbf{Use cases}: per-frame scratch memory, loading screens, temporary
computation buffers. The pattern is: \emph{allocate everything, process,
reset at frame end}. No individual deallocations are ever needed.

\begin{figure}[H]
\centering
\begin{tikzpicture}[scale=0.9]
  \draw[thick] (0,0) rectangle (8,0.6);
  \fill[darkblue!20] (0,0) rectangle (4.5,0.6);
  \draw[darkblue, thick, ->] (4.5,0.8) -- (4.5,0.6) node[above, yshift=4pt]{\small cursor};
  \node at (2.25,0.3) {\small used};
  \node at (6.25,0.3) {\small free};
\end{tikzpicture}
\caption{Linear allocator state: the cursor divides the buffer into used and free regions.}
\end{figure}

\section{Pool Allocator}

\begin{definition}[Pool Allocator]
A pool allocator divides a buffer into $N$ fixed-size \emph{slots} of
size $s$. Free slots are linked in a \emph{free-list}: each free slot
stores a pointer to the next free slot. Allocation pops the head of the
list; deallocation pushes back onto the head.
\end{definition}

\textbf{Complexity}: $O(1)$ allocation, $O(1)$ deallocation.
The list is embedded in-place: no separate bookkeeping array.

The free-list is initialised in reverse order so that the first
allocation returns the slot at the lowest address:

\begin{lstlisting}[caption={Free-list initialisation (reversed)}]
for (usize i = 0; i < m_maxSlots; ++i) {
    u8* slot = m_poolStart + i * m_slotSize;
    *reinterpret_cast<u8**>(slot) = m_freeList;
    m_freeList = slot;
}
\end{lstlisting}

\textbf{Use cases}: ECS component pools, audio source instances,
particle systems --- any subsystem that creates and destroys many
identically-sized objects at high frequency.


\paragraph{Invariant 4.2 --- Slot size}

  No allocation request may exceed \texttt{m\_slotSize}.
  The pool does not track individual allocation sizes; a larger
  request would corrupt adjacent slots.



\section{Stack Allocator}

The stack allocator is structurally identical to the linear allocator
but adds the concept of a \emph{Marker}:

\begin{lstlisting}[caption={Marker-based rollback}]
Marker m = stack.setMarker();    // snapshot cursor
// ... temporaries ...
stack.freeToMarker(m);           // roll back to m
\end{lstlisting}

A marker is simply a byte offset from the buffer start (\texttt{usize}).
\texttt{freeToMarker} moves the cursor back to the saved offset,
effectively deallocating everything allocated since the marker was set.
This supports LIFO (last-in, first-out) deallocation at $O(1)$ cost.

\textbf{Use cases}: recursive descent parsing, multi-phase asset loading
where each phase's scratch memory must be freed before the next phase.
